% Part: normal-modal-logic
% Chapter: frame-correspondence
% Section: introduction

\documentclass[../../../include/open-logic-section]{subfiles}

\begin{document}

\olfileid{nml}{frd}{int}

\olsection{పరిచయం}

మోడల్ తార్కికులకు ఆసక్తికరమైన ఒక ప్రశ్న:
ప్రాప్యత సంబంధానికీ, ఆ సంబంధం గల నమూనాల్లో కొన్ని
\tetoken{సూత్రాల}{!!{formula}s} సత్యానికీ మధ్య సంబంధం ఏమిటి?
ఉదాహరణకు, ప్రాప్యత సంబంధం స్వావర్తనమైనదని,
అంటే ప్రతి $w \in W$కు $Rww$ అని అనుకుందాం.
మరోలా అంటే, ప్రతి లోకం తనకు తానే ప్రాప్యమవుతుంది.
అప్పుడు ఒక లోకం~$w$ వద్ద $\Box !A$ సత్యమైతే,
$!A$ సత్యమై ఉండాల్సిన ప్రాప్య లోకాల్లో $w$
కూడా ఉంటుంది. కనుక $\mModel{M}$ నమూనాలోని
ప్రాప్యత సంబంధం~$R$ స్వావర్తనమైతే, ఏ లోకం $w$ను,
ఏ సూత్రం $!A$ను తీసుకున్నా అక్కడ $\Box !A
\lif !A$ సత్యమవుతుంది (మరోలా అంటే, $\Box p \lif
p$ పథకం, దాని ప్రతిస్థాపన నిదర్శనాలన్నీ
$\mModel{M}$లో సత్యమవుతాయి).

అయితే, దీనికి విలోమం సరికాదు. ఉదాహరణకు,
$\Box p \lif p$ అనేది $\mModel{M}$లో సత్యమైతే
$R$ స్వావర్తనమని తేలదు. ఎందుకంటే స్వావర్తనం కాని,
అయినా ప్రతి లోకంలో $\Box p \lif p$ సత్యమయ్యే
నమూనా~$\mModel{M}$ను సులభంగా చూపవచ్చు:
ఒకే లోకం~$w$ ఉండి, అది తనకు తానే ప్రాప్యం కాని,
$w \in V(p)$ అయ్యే నమూనాను తీసుకుందాం.
ఈ ఉదాహరణలో $!A=p$గా తీసుకొని $p$ సత్యమూల్యాన్ని
తగినట్లు ఎంచుకుంటే, స్వావర్తనం కాని నమూనాలోనూ
$\Box !A \lif !A$ సత్యమయ్యేలా చేయవచ్చు.
\sourcecorrection{OLTENMLFRDINT-001}{మూలం $p$కు
  సత్యమూల్యం ఎంచుకున్న ఉదాహరణ నుంచి సాధారణ
  $!A$ రూపానికి మారుతుంది. ఇక్కడ $!A=p$ సందర్భానికే
  ఆ నిర్ధారణను పరిమితం చేశాం; ఏకపక్ష $!A$కు
  నిరూపణగా పరిగణించలేదు.}

పరిష్కారం: చర కేటాయింపు~$V$ను ఈ పరిశీలన నుంచి
తొలగించాలి. $V(p)$లో ఏ లోకాలు ఉన్నా,
$\mModel{M}$లోని అన్ని లోకాల్లో $\Box p \lif p$
సత్యం కావాలని కోరితే, $R$ తప్పక స్వావర్తనమై ఉండాలి.
ఎందుకంటే స్వావర్తనం కాని ఏ నమూనాలోనైనా
$Rww$ కాని లోకం~$w$ కనీసం ఒకటి ఉంటుంది.
$V(p) = W \setminus \{w\}$గా తీసుకుంటే,
$p$ అనేది $w$ తప్ప మిగిలిన అన్ని లోకాల్లో సత్యం;
కాబట్టి $w$ నుంచి ప్రాప్యమైన అన్ని లోకాల్లోనూ
సత్యం ($w$ నుంచి $w$కు ప్రాప్యం కాదనీ,
$p$ అసత్యమయ్యే ఏకైక లోకం $w$ అనీ తెలుసు).
కానీ $w$ వద్ద $p$ అసత్యం కాబట్టి
$\Box p \lif p$ సూత్రం $w$ వద్ద అసత్యం.

అందువల్ల మూల్యనిర్ణయం లేని నమూనా నిర్మాణాలకు
ఒక సంకేతం అవసరమని తెలుస్తుంది: వాటిని
\emph{చట్రాలు} అంటాం. చట్రం $\mModel{F}$ అంటే
ప్రాప్యత సంబంధంతో కూడిన లోకాల సమితి గల
జత $\tuple{W, R}$ మాత్రమే. అప్పుడు ప్రతి నమూనా
$\tuple{W, R, V}$ ఆ చట్రం $\tuple{W, R}$పై
\emph{ఆధారపడుతుంది} అంటాం. విలోమంగా,
ఒక చట్రం దానిపై ఆధారపడే నమూనాల వర్గాన్ని
నిర్ణయిస్తుంది; చట్రాల వర్గం ఆ వర్గంలోని ఏదో
చట్రంపై ఆధారపడే నమూనాల వర్గాన్ని నిర్ణయిస్తుంది.
ఒక చట్రంలో \tetoken{సూత్రం}{!!a{formula}}
\emph{చెల్లుబాటు} అవుతుందనే భావాన్ని
$\mModel{F} \Entails !A$గా నిర్వచించవచ్చు:
$\mModel{F}$పై ఆధారపడే అన్ని $\mModel{M}$
నమూనాలకు $\mSat{M}{!A}$ ఉండాలి.

ఈ సంకేతంతో \tetoken{సూత్రాల}{!!{formula}s}కీ
చట్రాల వర్గాలకీ మధ్య అనురూపతా సంబంధాలను
స్థాపించవచ్చు: ఉదాహరణకు,
$\mModel{F} \Entails \Box p \lif p$
అప్పుడే, అప్పుడే $\mModel{F}$ స్వావర్తనమైనది.

\end{document}
